% =====================================================================
% Chapter 3: Methodology -- revised draft
% Required packages (in addition to your existing preamble):
%   amsmath, amssymb, graphicx, float, natbib, todonotes, array, url
% References for citations in this file are included below; no .bib file is needed.
% Items marked \todo{CONFIRM ...} require information that is not
% recorded in the project repository and must be verified by the author.
% =====================================================================

\section{Methodology}

\subsection{Introduction}

This chapter describes the methodology used to investigate curriculum learning
for automatic speech recognition (ASR) in Shona (ChiShona), a Bantu language
spoken mainly in Zimbabwe for which comparatively little transcribed speech is
available. Curriculum learning is a training strategy in which examples are
presented to a model in a deliberately chosen order, typically from examples
judged to be easier towards examples judged to be harder, rather than in a
uniformly random order \citep{Bengio2009}. We adopted an experimental approach
to determine whether ordering training data according to an estimate of sample
difficulty improves the recognition accuracy of a pretrained ASR model that is
fine-tuned on limited Shona speech. Fine-tuning refers to continuing the
training of a model that has already been trained on a large general corpus,
using a smaller task-specific dataset, so that the general knowledge of the
model is adapted to the new language or domain.

The chapter first describes the research design and the computing environment.
It then describes how the speech data were acquired, inspected, cleaned,
normalized, and re-partitioned so that the speakers used for evaluation never
appear in training. Next, it describes the baseline ASR model, the curriculum
learning strategies, and the data augmentation techniques. Data augmentation
refers to artificially modifying training examples, for example by adding
background noise, so that the model is exposed to more acoustic variability than
the original recordings contain. The chapter then describes a model-capacity
extension in which the principal experiments were repeated with a larger model.
Finally, it defines the evaluation metrics and the statistical procedures used
to compare training strategies.

Word error rate (WER) was the primary evaluation metric and character error rate
(CER) was the secondary metric; both are defined formally in
Section~\ref{sec:evaluation}. Every training strategy was compared with a
matched baseline that differed from it in only one factor, so that any
difference in WER could be attributed to that factor.

\subsection{Research Design}
\label{sec:research-design}

\subsubsection{Experimental approach}

We approached this study with a quantitative experimental design. ASR systems
were trained under controlled conditions and evaluated on utterances that were
not seen during training. Evaluation on unseen utterances was necessary because
the purpose of an ASR system is to transcribe new speech; performance on the
training utterances would measure how well the model had memorized those
specific recordings rather than how well it generalizes. For the same reason, the
principal evaluation partitions were constructed so that their speakers were
also absent from training (Section~\ref{sec:speaker-disjoint}).

\subsubsection{Variables}

The \textbf{primary independent variable} was the curriculum schedule, that is,
the order in which the labelled training utterances were presented to the model
during fine-tuning. Three \textbf{secondary independent variables} were studied:

\begin{enumerate}
    \item the \emph{data augmentation policy}, which was either no augmentation,
    synthetic waveform augmentation (additive noise, simulated reverberation, and
    tempo change applied to the audio signal), or spectrogram masking
    augmentation (SpecAugment, applied to the model's input features);
    \item the \emph{interaction} between the curriculum schedule and the
    augmentation policy, examined through factorial designs in which every
    combination of the two factors was trained; and
    \item the \emph{model capacity}, examined by repeating the principal
    conditions with Whisper Medium (769 million parameters) in addition to
    Whisper Base (74 million parameters).
\end{enumerate}

The \textbf{primary dependent variable} was corpus-level WER, because WER is the
standard measure of ASR transcription accuracy and allows our results to be
compared with published work. \textbf{Secondary dependent variables} were CER and
the evaluation loss. CER was included because Shona is an agglutinative language
in which words are formed by joining several meaningful units (morphemes); a
single incorrect morpheme makes the whole word incorrect under WER, so CER can
reveal partial improvements that WER conceals. The evaluation loss (the
cross-entropy of the reference transcription under the model, defined in
Section~\ref{sec:input-preprocessing}) was recorded as an optimization diagnostic
only. It was not used to judge transcription quality, because a model can assign
higher probability to the reference text without producing a better transcription
when it decodes freely.

The \textbf{controlled variables} were the pretrained model and its exact
version, the data partitions, the text normalization, the number of training
epochs, the optimizer and learning-rate schedule, the batch size, the random seed,
the decoding configuration, and the checkpoint-selection rule. These were held
constant within every comparison so that only the factor under investigation
differed.

\subsubsection{Research questions}

The methodology was designed to answer the following research questions:

\begin{enumerate}
    \item[\textbf{RQ1:}] What are the effects of static and dynamic acoustic
    data-level curriculum learning on the fine-tuning of ASR foundation models
    on limited Shona speech?
    \item[\textbf{RQ2:}] How does combining curriculum schedulers with synthetic
    data augmentation affect the in-domain and out-of-domain generalization of
    Shona ASR models?
\end{enumerate}

An ASR \emph{foundation model} is a large model pretrained on very large and
diverse speech corpora that can be adapted to many downstream tasks; Whisper
\citep{Radford2022} is the foundation model used in this study. \emph{Static}
curricula compute the difficulty of each utterance once, before training, from
properties of the recording; \emph{dynamic} curricula update difficulty during
training from the behaviour of the model itself (Section~\ref{sec:curriculum}).
\emph{In-domain} generalization refers to unseen speakers recorded under the same
collection protocol as the training data, while \emph{out-of-domain}
generalization refers to speech collected through a different process, with
different speakers, recording conditions, and text topics
(Section~\ref{sec:fleurs}).

\subsubsection{Experimental sequence}

The experiments were carried out in the following order. Each step depended on
the outcome of the previous one, and the order was fixed before results were
inspected so that later choices could not be tuned to favour a preferred
conclusion.

\begin{enumerate}
    \item Construct and freeze a speaker-disjoint data protocol, the transcript
    normalization, and the acoustic measurements used for difficulty scoring.
    \item Establish a baseline that used conventional random ordering and no
    augmentation.
    \item Compare static, dynamic, and hybrid curriculum strategies against that
    baseline with Whisper Base (RQ1).
    \item Evaluate data augmentation independently of curriculum, and in
    combination with the best-performing curriculum, using factorial designs
    (RQ2).
    \item Evaluate the in-domain models on an out-of-domain corpus only after the
    curriculum choice had been fixed using in-domain validation data (RQ2).
    \item Repeat the principal conditions with Whisper Medium to determine whether
    the conclusions depend on model capacity (Section~\ref{sec:medium}).
\end{enumerate}

\begin{figure}[H]
    \centering
    \includegraphics[width=1\linewidth]{pipeline_diagram.png}
    \caption{Overview of the experimental pipeline, from data acquisition and
    speaker-disjoint partitioning to curriculum and augmentation experiments,
    model-capacity extension, and evaluation.}
    \label{fig:pipeline}
\end{figure}

\subsubsection{Exploratory and controlled phases}

The work was conducted in two phases. The \textbf{exploratory phase} used
interactive notebooks to investigate model sizes, preliminary curriculum
strategies, and preliminary augmentation settings, and to build the data
preparation and training pipeline. Because these exploratory runs often changed
more than one factor at a time, and because they used the original WAXAL
partitions in which most evaluation speakers also appear in training
(Section~\ref{sec:leakage}), they were retained only as exploratory evidence that
informed later design choices. They were not used to make causal claims.

The \textbf{controlled phase} used a configuration-driven pipeline in which every
experiment was defined by a versioned configuration file, every data file was
identified by a cryptographic hash, the exact version of the pretrained model was
fixed, and random seeds were set explicitly. All conclusions reported for RQ1 and
RQ2 are drawn from this phase.

\subsubsection{Controlled comparisons, random seeds, and replication}
\label{sec:seeds}

We treated a comparison as controlled only when the dataset, the training budget,
the random seed, the decoding configuration, and the evaluation partitions were
identical between the two conditions being compared. Earlier runs that altered
more than one of these factors were retained as exploratory evidence.

Training a neural network involves several random processes, including the order
in which examples are shuffled, the initialization of any newly created
parameters, and stochastic augmentation decisions. A random seed is the number
used to initialize the pseudo-random number generator, so the same seed
reproduces the same sequence of random decisions. Different seeds can produce
noticeably different final models even when every other setting is identical
\citep{Reimers2017, Dodge2020}. For this reason the controlled experiments were
repeated with three seeds (42, 43, and 44). The same three seeds were shared
across conditions, so that each treatment run could be paired with a control run
that experienced the same random decisions wherever the treatment allowed. Three
seeds were chosen as the largest number that the available computing budget
permitted across all conditions; this allowed us to observe run-to-run
variability and to avoid attributing a lucky or unlucky seed to the treatment.
Three seeds are, however, too few to estimate the full distribution of seed
effects, and this limitation is considered when results are interpreted.

Not every condition was replicated from the outset. The initial curriculum
screen trained each condition with seed 42 only, in order to remove clearly
unpromising strategies cheaply. Seeds 43 and 44 were added afterwards for the
baseline and for selected conditions (C1, C2, C3, C5, and C6). Because these
additional runs were chosen after the seed-42 results had been inspected, they
are reported as post hoc replications rather than as part of a single
prespecified three-seed trial.

Two curricula (C2 and C3, Section~\ref{sec:curriculum}) present the training data
in a fully deterministic order that does not depend on the seed, and Whisper's
dropout probability is zero in its released configuration. These two conditions
therefore produced byte-identical models for all three seeds. Their repeated runs
demonstrate exact reproducibility rather than independent replication, and their
results are compared against the three independently seeded baseline runs to
quantify the variability of the control.

\subsection{Implementation Environment}
\label{sec:implementation}

The training pipeline was implemented in Python 3.10.12 using PyTorch 2.5.1
\citep{Paszke2019}, an open-source deep-learning library. The following
additional packages were used, each with the role described:

\begin{itemize}
    \item \textbf{Transformers 4.46.3} \citep{Wolf2020}: provides the pretrained
    Whisper models, the feature extractor, the tokenizer, and the
    \texttt{Seq2SeqTrainer} training loop;
    \item \textbf{Datasets 2.20.0} \citep{Lhoest2021}: loads the data manifests and
    decodes and resamples audio on demand;
    \item \textbf{Accelerate 1.1.1}: manages device placement and mixed-precision
    training;
    \item \textbf{Evaluate 0.4.2} and \textbf{JiWER 3.0.4}: compute WER and CER
    (Section~\ref{sec:metric-computation});
    \item \textbf{NumPy 1.26.4}: numerical computation for audio analysis and
    statistics;
    \item \textbf{SoundFile}, \textbf{librosa} \citep{McFee2015}, and
    \textbf{SciPy}: reading and writing audio files, pitch-preserving tempo
    modification, and convolution for waveform augmentation;
    \item \textbf{bitsandbytes 0.45.5} \citep{Dettmers2022}: provides an 8-bit
    optimizer that was used only in the exploratory Whisper Large run, not in
    any controlled run; and
    \item \textbf{Weights \& Biases (W\&B) 0.28.2} \citep{Biewald2020}: records
    training curves, configurations, and system metrics for every run.
\end{itemize}

All controlled experiments were trained on an NVIDIA GeForce RTX 4090 graphics
processing unit (GPU) with 24~GB of video memory (VRAM), using CUDA 12.4, NVIDIA's
programming platform for GPU computation.
\todo{CONFIRM: an NVIDIA A100 (40~GB) appears only in a legacy cluster
scheduler script; no controlled-run manifest records an A100. If the A100 was
used, state which exploratory runs used it and why (for example, runs that did
not fit in 24~GB). Otherwise remove it.}
Using a single GPU type for all controlled comparisons was a deliberate decision:
different GPU architectures perform floating-point arithmetic in slightly
different orders, which can change results, so training a treatment and its
control on different hardware would introduce an uncontrolled factor.

Each experiment was defined by a versioned configuration file in JavaScript
Object Notation (JSON), a plain-text format for structured data. Each
configuration recorded the model identifier and its exact repository revision
(a fixed version identifier that prevents the model from silently changing if
its publisher updates it), the paths of the data manifests, the language and
task, the random seed, the batch and gradient-accumulation settings, the
optimization schedule, the augmentation policy where one was applied, the output
location, and the experiment-tracking metadata. Before training, the pipeline
wrote a run manifest containing the Secure Hash Algorithm 256-bit (SHA-256) hash
of the training script, the configuration, the dependency list, and every data
manifest, together with the package versions, the GPU model, and the effective
seed. A SHA-256 hash is a 64-character fingerprint computed from the contents of a
file; any change to the file, however small, produces a different hash.
Recording hashes therefore made it possible to prove, after the fact, that two
runs used exactly the same data and code. Long-running jobs were executed inside
detached terminal sessions so that training continued if the remote connection
to the server was interrupted, and resource checks confirmed that sufficient GPU
memory and disk space were available before each run was launched.

\subsection{Dataset and Data Preparation}
\label{sec:data}

\subsubsection{The WAXAL corpus}

This study used the Shona component of the WAXAL speech corpus \citep{Diack2026}
as the primary source of training and evaluation data. WAXAL was collected
through a multi-year effort funded by Google in partnership with organizations
across the African continent. The Shona component contains 99.2 hours of
transcribed speech and a total of 574.2 hours of speech when untranscribed
recordings are included. Only the transcribed portion was used in this study.
The collection process was designed to elicit natural rather than read speech
and involved the following stages:

\begin{itemize}
    \item \textbf{Image-prompted speech:} participants were shown images covering
    50 topics and asked to describe what they saw in their native language.
    Describing an image, rather than reading a prepared sentence, produces
    spontaneous speech with natural hesitations and vocabulary.
    \item \textbf{Recording:} recordings were captured in the speakers' natural
    environments rather than in a studio, so they contain realistic background
    noise and room acoustics. \citet{Diack2026} describe a minimum recording
    length of 15 seconds as part of the collection protocol; the released Shona
    clips used in this study are nevertheless shorter in some cases (3.6 to 35.9
    seconds after preparation), so the 15-second figure describes the provider's
    elicitation protocol rather than every released clip.
    \item \textbf{Transcription:} approximately 10\% of the audio collected
    across the whole WAXAL project was transcribed by paid professional local
    linguists \citep{Diack2026}. This proportion refers to the overall
    multilingual collection; it should not be computed from the Shona hours
    alone. Each recording was transcribed by a single annotator.
    Because transcriptions were not cross-checked by a second annotator, some
    transcription errors are expected; Section~\ref{sec:filtering} describes the
    automatic checks used to detect the most severe mismatches.
    \item \textbf{Quality control:} each recording was checked for clarity,
    language accuracy, and relevance to the prompt, and personally identifiable
    information was removed where present.
\end{itemize}

Each released record contains the audio, a reference transcription, an
utterance identifier, a speaker identifier, the language, and, where available,
the speaker's gender. The speaker identifier is essential to this study because
it makes it possible to construct evaluation partitions whose speakers are absent
from training (Section~\ref{sec:speaker-disjoint}).

\subsubsection{Acquisition and manifest conversion}
\label{sec:acquisition}

Audio and transcript analysis was conducted with the NVIDIA NeMo Speech Data
Explorer (SDE) \citep{Bakhturina2021}. NeMo is NVIDIA's open-source toolkit for
speech and language models, and SDE is an interactive, browser-based tool within
it that displays a speech dataset as sortable tables and plots and allows any
utterance to be played back alongside its transcription, waveform, and
spectrogram. We chose SDE because it was designed specifically for auditing ASR
datasets: it computes duration, character-rate, and word-rate distributions
automatically, highlights out-of-vocabulary characters, and allows outliers to be
listened to directly, without writing custom visualization code.

SDE reads data in the NeMo manifest format, which is a JSON Lines (JSONL) file:
a plain-text file in which each line is one JSON object describing one
utterance. We therefore wrote a script that loaded the published training,
validation, and test partitions of the WAXAL Shona ASR subset from the Hugging
Face dataset repository \texttt{google/WaxalNLP}, skipped records whose
transcription was empty, and wrote one manifest per partition. Each audio clip was
saved as a separate local file, organized by partition, in its original encoding;
the audio was not re-encoded, because re-encoding to a lossy format could
introduce compression artefacts. The duration of each clip was computed from the
audio file itself, as the number of audio samples (frames) divided by the
sampling rate, using SoundFile, a Python library that reads audio files through
the libsndfile C library. We measured duration directly rather than trusting the
duration recorded in the dataset metadata because later steps (rate-based
filtering, silence trimming, and duration-based curricula) depend on duration
being exact. Each manifest row stored the absolute audio path, the measured
duration, the transcription, the utterance identifier, and, where available, the
speaker, language, and gender fields.

The resulting manifests contained 14,109 training, 1,727 validation, and 1,749
test utterances.

\subsubsection{Quality analysis and rule-based filtering}
\label{sec:filtering}

We used SDE to inspect utterance duration, transcription length, speaking-rate
outliers, unexpected characters, and other anomalies. The patterns observed
during this inspection were then converted into fixed rules that were applied by
a script. This division of labour was deliberate: SDE supported exploration and
listening, while the actual decisions to keep or reject an utterance were made by
code, so that they were reproducible and documented rather than dependent on
manual judgement.

For each utterance $i$ we calculated the character rate and the word rate.
These rates measure how much text is associated with each second of audio. For
correctly transcribed speech they fall within a predictable range, so an extreme
value indicates that the transcription and the audio do not match, for example
because the transcription belongs to a different recording, describes only part
of the recording, or the audio was truncated:

\begin{equation}
    r_i^{\text{char}} = \frac{|y_i|}{d_i}, \qquad
    r_i^{\text{word}} = \frac{\operatorname{words}(y_i)}{d_i},
    \label{eq:rates}
\end{equation}

where
$r_i^{\text{char}}$ is the character rate (characters per second) of utterance $i$;
$r_i^{\text{word}}$ is the word rate (words per second) of utterance $i$;
$i$ is the index identifying a specific audio clip and its transcription;
$y_i$ is the reference transcription of utterance $i$;
$|y_i|$ is the number of characters in $y_i$, including spaces and any
punctuation present in the supplied transcription;
$\operatorname{words}(y_i)$ is the number of whitespace-separated words in $y_i$; and
$d_i$ is the duration of the audio clip in seconds.

An utterance was automatically rejected if any of the following conditions held:

\begin{enumerate}
    \item \textbf{Duration below one second.} The WAXAL elicitation task asked
    speakers to describe an image, which cannot be done meaningfully in less than
    a second. A clip shorter than one second therefore almost certainly
    represents a failed, truncated, or accidental recording rather than a valid
    utterance.
    \item \textbf{Character rate below 1 or above 25 characters per second.}
    Conversational speech typically corresponds to roughly 10--15 written
    characters per second. A rate below 1 means that most of the audio has no
    corresponding text (for example, a transcription covering only the first few
    words), while a rate above 25 means that more text is attributed to the audio
    than could plausibly have been spoken in that time. Both thresholds are
    deliberately wide, so that only clear mismatches are removed and naturally
    fast or slow speakers are retained.
    \item \textbf{Word rate above 4 words per second.} Typical conversational
    speaking rates are approximately 2--3 words per second \citep{Yuan2006}, and
    Shona words are, on average, long because of agglutination, which lowers the
    expected word rate further. A rate above 4 words per second is therefore
    implausible for genuine Shona speech and indicates a transcription mismatch.
\end{enumerate}

The thresholds were applied to all three partitions, not only to the training
data. A mismatched reference in a validation or test partition would penalize a
model for failing to produce words that were never spoken, inflating WER with
errors that reflect the dataset rather than the model. Filtering the evaluation
partitions with the same rules therefore made the evaluation a more valid
measure of recognition accuracy, and applying identical rules everywhere kept the
three partitions statistically comparable. Because the same rules were fixed
before any model was trained, they could not have been tuned to favour a
particular model. The thresholds are local screening choices informed by the
inspection described above. They are not validated universal boundaries for
Shona, and no separate threshold-optimization study was carried out. The
preserved rejection manifests show that every one of the 392 rejected
utterances violated both the upper character-rate limit and the word-rate
limit; the duration rule and the lower character-rate rule did not, by
themselves, reject any utterance in this export. In practice, therefore, the
filter removed utterances whose transcriptions contained far more text than
their audio could contain.

The script also flagged transcriptions containing embedded line breaks or
characters outside the expected inventory. Flagged utterances were reviewed
manually rather than rejected automatically, because some characters that appear
unusual are linguistically meaningful in Shona. For example, the apostrophe in
\emph{n'anga} (traditional healer) marks the velar nasal sound, written
\emph{n'} and pronounced like the \emph{ng} in English \emph{sing}, which
distinguishes it from the sequence \emph{n} + \emph{g}. Removing such characters
automatically would change the words and their pronunciation.

Table~\ref{tab:filtering} summarizes the outcome. In total, 392 utterances
(2.2\%) were rejected.

\begin{table}[H]
    \centering
    \begin{tabular}{|l|r|r|r|}
        \hline
        Partition & Input & Retained & Rejected \\
        \hline
        Training & 14,109 & 13,799 & 310 \\
        \hline
        Validation & 1,727 & 1,683 & 44 \\
        \hline
        Test & 1,749 & 1,711 & 38 \\
        \hline
    \end{tabular}
    \caption{Number of utterances retained and rejected by the rule-based quality
    filter in each original WAXAL partition.}
    \label{tab:filtering}
\end{table}

For each partition the cleaning process produced five artefacts: (i) a cleaned
manifest containing the retained utterances; (ii) a rejected manifest listing
every rejected utterance together with the rule that caused its rejection; (iii)
an enriched manifest containing all derived measurements; (iv) a
comma-separated values (CSV) file for manual review in a spreadsheet; and (v) a
JSON policy record. The policy record stored the fixed thresholds, the SHA-256
hash of the input manifest, the numbers of retained and rejected utterances, and
the output paths. No source audio was deleted, so every decision can be reversed
or audited.

\subsubsection{Silence trimming and acoustic audit}
\label{sec:silence}

We examined the retained training audio for long periods of silence. Long
silences are undesirable for two reasons. First, Whisper processes audio in
fixed 30-second windows (Section~\ref{sec:input-preprocessing}), so silence
occupies part of the window without contributing information. Second, Whisper is
known to produce fluent but entirely fabricated text, often called
\emph{hallucination}, when given long stretches of non-speech audio
\citep{Koenecke2024}; training on silence paired with text could reinforce this
behaviour. Long silences also distort the duration-based measurements used for
filtering and for curriculum scoring.

Each recording was divided into consecutive 100-millisecond frames, and the
root mean square (RMS) energy of each frame was computed:

\begin{equation}
    \operatorname{RMS}_f = \sqrt{\frac{1}{M}\sum_{m=1}^{M} x_{f,m}^2},
    \label{eq:rms}
\end{equation}

where $x_{f,m}$ is the $m$-th audio sample in frame $f$ and $M$ is the number of
samples per frame. A frame was classified as \emph{active} (containing speech or
other sound) when its RMS energy was within 40 decibels (dB) of the loudest frame
in the same recording:

\begin{equation}
    20\log_{10}\!\left(\frac{\operatorname{RMS}_f}{\max_{g}\operatorname{RMS}_g}\right) \geq -40~\text{dB}.
    \label{eq:active}
\end{equation}

A threshold of $-40$~dB corresponds to an amplitude of 1\% of the loudest frame.
The threshold is defined relative to each recording's own peak rather than as an
absolute level because WAXAL recordings were captured on different devices at
different recording levels; a relative threshold adapts automatically to each
device. The value of 40~dB is a study-specific heuristic: it is low enough that
quiet speech sounds such as word-final consonants remain classified as active,
yet high enough that the steady background hiss of a typical consumer recording
is classified as inactive. The same definition was later reused for the
signal-to-noise ratio estimate (Section~\ref{sec:snr}), so that all acoustic
measurements in the study share one definition of activity.

When at least two seconds of trailing silence were detected after the last active
frame, the recording was trimmed after that frame, retaining 250 milliseconds of
padding. The two-second threshold ensured that only clearly excessive tails were
removed, while the natural short pauses that occur at the end of spoken
utterances were left untouched. The 250-millisecond padding protected the decaying
end of the final word, because low-energy sounds such as fricatives and voice
decay can fall below the activity threshold even though they belong to the
speech. Leading silence was measured and flagged for review but was not trimmed,
because the start of a recording was more likely than the end to contain quiet
speech onsets.

The trimmed audio was saved as a new file in the Free Lossless Audio Codec (FLAC)
format, and the original file was preserved. FLAC was chosen because it is
lossless, meaning that the decoded audio is sample-for-sample identical to the
input, while occupying substantially less disk space than uncompressed audio; it
is also read natively by SoundFile and by the Hugging Face audio loader. Trimming
was applied only to the original WAXAL training partition, where it affected 530
of the 13,799 retained recordings; the remaining 13,269 recordings were
unchanged, and the original validation and test recordings were not trimmed.
When the three partitions were later pooled and re-divided by speaker
(Section~\ref{sec:speaker-disjoint}), the 530 trimmed recordings were
distributed across the new partitions: 434 entered the new training partition,
58 the new validation partition, and 38 the new test partition. Consequently, a
small number of held-out recordings (5.7\% of new test utterances or fewer)
had undergone label-independent trimming while the others had not. Trimming
removes only trailing non-speech and does not use the transcription, so it
cannot leak label information, but the asymmetry is disclosed here because it
means the held-out partitions were not processed entirely uniformly. A CSV audit file linked every original path and
duration to its output path, its new duration, and the number of seconds removed.

An acoustic audit then decoded every retained training recording and measured
duration agreement with the manifest, sampling rate, number of channels, peak and
average RMS level, clipping (the fraction of samples at the maximum representable
amplitude, which indicates distortion), direct-current offset (a constant shift of
the waveform away from zero, which indicates a faulty recording chain), the ratio
of active to total frames, leading and trailing silence, and a signal-to-noise
ratio estimate. All training files decoded successfully and matched their manifest
durations, and no clipping or residual long trailing silence was found.
Recordings that were potentially difficult but valid, such as those with long
leading silence, a low proportion of active speech, or low amplitude, were flagged
for review rather than removed, because removing them would make the training data
artificially clean and less representative of real-world speech.

\subsubsection{Transcript normalization}
\label{sec:normalization}

WER and CER are computed by comparing the model's output with the reference text
character by character and word by word. Any difference in the written form of
the same spoken words, such as capitalization or punctuation, would therefore be
counted as a recognition error even though it carries no acoustic information.
To prevent such spurious errors, we applied a single text-normalization procedure
(the \emph{text contract}) to the training, validation, and test partitions, and
the same procedure was later applied to the out-of-domain corpus
(Section~\ref{sec:fleurs}). The procedure consisted of the following steps.

\begin{enumerate}
    \item \textbf{Unicode NFC normalization.} Unicode, the international standard
    for encoding text, allows some characters to be encoded in more than one way;
    for example, an accented letter can be stored either as a single precomposed
    character or as a base letter followed by a separate combining accent. The two
    forms look identical but are different sequences of codes, so a string
    comparison would treat them as different. Normalization Form Canonical
    Composition (NFC) \citep{UAX15} converts every such character to a single
    canonical form, ensuring that visually identical text is also identical for
    comparison.
    \item \textbf{Apostrophes and dashes.} Typographic (curly) apostrophes
    (\texttt{'} and similar characters) were mapped to the plain American Standard
    Code for Information Interchange (ASCII) apostrophe, and the various
    typographic dashes were standardized. Different transcribers and keyboards
    produce different apostrophe characters for the same Shona sound; without this
    step, \emph{n'anga} typed with a curly apostrophe and with a straight
    apostrophe would be counted as different words.
    \item \textbf{Case folding.} All text was converted to lower case. Capital
    letters are a writing convention that cannot be heard in speech, and the
    WAXAL transcriptions applied capitalization inconsistently, so retaining case
    would penalize the model for orthographic choices rather than recognition
    errors.
    \item \textbf{Punctuation removal.} Letters, digits, and apostrophes were
    retained; all other punctuation and symbols were replaced by spaces, and
    repeated spaces were then collapsed into a single space. Punctuation is not
    spoken, and transcribers punctuated inconsistently. Punctuation was replaced
    with a space rather than deleted so that two words separated only by a
    punctuation mark (for example, \emph{word,word}) were not merged into a single
    incorrect word. Apostrophes were retained because, as explained in
    Section~\ref{sec:filtering}, they are part of Shona spelling.
\end{enumerate}

The original transcription was preserved in a separate \texttt{text\_raw} field
and the normalized form was stored in the \texttt{text} field. Retaining the
original made the normalization fully reversible and auditable: any normalized
transcription can be compared with its source to verify that the procedure did
not remove meaningful content, and a different normalization could be applied in
future work without re-downloading the corpus. Normalization produced no empty
transcriptions and no transcriptions in a non-Latin script.

Twenty-seven training and four test transcriptions in the original partitions
contained digits (for example, \emph{5} rather than the spoken Shona number
word). These were flagged for manual review
rather than converted automatically. Speakers say number words, not digits, so a
digit in the reference cannot match any spoken form; however, automatic
conversion would require a Shona-specific number-to-word converter, and applying
English conventions would produce incorrect Shona. Because the affected
utterances represent approximately 0.2\% of the data, they were retained
unchanged and documented as a known source of minor label noise.

\subsubsection{Duplicate and split-leakage audit}
\label{sec:leakage}

We checked whether any recording occurred more than once, either within a
partition or across partitions. If the same audio appeared in both the training
and test partitions, the model could recognize it from memory, inflating test
accuracy. This type of contamination is known as \emph{data leakage}
\citep{Kaufman2012}.

To detect duplicates, every audio file was decoded into its raw pulse-code
modulation (PCM) samples, the sequence of numerical amplitude values that
represents the digital waveform, and a SHA-256 hash was computed over those
samples. Hashing the decoded samples rather than the file itself was necessary
because two files can contain identical audio but differ in their headers,
metadata, or compression settings, which would give them different file hashes;
two recordings with identical decoded samples, by contrast, always produce the
same PCM hash. No exact audio duplicates were found within any partition or
across the training, validation, and test partitions. Two groups of repeated
transcriptions, comprising four training records, were retained because their
audio differed, which indicates different speakers or different takes of the same
description. One normalized transcription was shared between the training and
validation partitions, but its audio was not duplicated.

The audit revealed a more serious form of overlap. The original WAXAL partitions
do not provide a speaker-independent test set: 94 of the 97 validation speakers
and 101 of the 110 test speakers also appear in the training partition, and 1,700
of the 1,711 retained test utterances were produced by speakers who are present in
training. A model evaluated on these partitions has already learned the voices,
accents, vocabulary, and recording devices of almost all the test speakers, so
the original test set measures recognition of familiar speakers more than
generalization to new ones. Because a deployed Shona ASR system would be used
predominantly by people whose voices were never seen during training, we
constructed new speaker-disjoint partitions for all controlled experiments, as
described in the next section. Standard ASR benchmarks such as TIMIT
\citep{Garofolo1993} and LibriSpeech \citep{Panayotov2015} use speaker-disjoint
evaluation sets for the same reason.

\subsubsection{Speaker-disjoint split construction}
\label{sec:speaker-disjoint}

\paragraph{Pooling and grouping.}
To construct the speaker-disjoint partitions, the three cleaned and normalized
WAXAL manifests were pooled into a single set of 17,193 utterances
($13{,}799 + 1{,}683 + 1{,}711$). Each utterance was then grouped by its speaker
identifier: a mapping was built from each of the 161 distinct speaker identifiers
to the list of all utterances recorded by that speaker. Assignment to the new
partitions was performed on these speaker groups rather than on individual
utterances, so that all utterances of a given speaker were placed in the same
partition and no speaker could appear in more than one partition. Every
utterance was required to have a speaker identifier; the script stopped with an
error if any identifier was missing.

\paragraph{Target proportions.}
The target proportions were 80\% of the data for training and 10\% each for
validation and test, matching the approximate proportions of the original WAXAL
release. Because speakers contribute very different numbers of utterances, it is
not possible to choose groups of whole speakers that reach these proportions
exactly; the task is therefore an optimization problem in which we seek the
assignment of speakers that comes closest to the targets.

\paragraph{Objective function.}
For a candidate assignment $a$ of speakers to partitions, the deviation from the
targets was measured by the objective

\begin{equation}
J(a) = \sum_{k \in \{\text{train}, \text{validation}, \text{test}\}} \left[ \left( \frac{n_k(a) - r_k N}{r_k N} \right)^2 + \left( \frac{h_k(a) - r_k H}{r_k H} \right)^2 \right],
\label{eq:split-objective}
\end{equation}

where $k$ indexes the three partitions; $r_k$ is the target proportion of
partition $k$ (0.8, 0.1, and 0.1); $n_k(a)$ is the number of utterances that
assignment $a$ places in partition $k$; $h_k(a)$ is the corresponding total
duration in hours; $N$ is the total number of utterances (17,193); and $H$ is the
total duration in hours. The quantities $r_k N$ and $r_k H$ are therefore the
target number of utterances and the target duration of partition $k$.

Each term in Equation~\eqref{eq:split-objective} is a \emph{squared relative
deviation}: the difference between the achieved and target size, divided by the
target size, and squared. Dividing by the target size makes each term a
proportion, so that a 5\% shortfall in the small test partition is penalized as
heavily as a 5\% shortfall in the large training partition; without this
normalization, the objective would be dominated by the training partition, and
the held-out partitions could deviate substantially from their targets. Squaring
makes all deviations positive and penalizes large deviations more than several
small ones, which favours assignments in which all partitions are close to their
targets. Both the number of utterances and the duration were included, with
equal weight, because they capture different properties: the number of utterances
determines how many independent items contribute to the evaluation, whereas the
duration determines how much speech the model is trained or evaluated on. Since a
few speakers recorded many short or many long utterances, balancing only one
measure could leave the other badly unbalanced. A perfect assignment would give
$J(a) = 0$. Because the three partition sizes must sum to the total, only two of
the three partitions vary independently for each measure; this does not affect
the procedure, because $J(a)$ is used only to rank candidate assignments, not as a
statistical test.

\paragraph{Concentration constraint.}
Minimizing $J(a)$ alone does not prevent a single speaker from dominating a
held-out partition. If one speaker contributed, for example, half of the test
utterances, the test WER would largely measure how well the model recognized that
one person. A candidate assignment was therefore eligible only if no single
speaker contributed more than 20\% of the utterances or more than 20\% of the
duration of the validation or test partition. A cap of 20\% guarantees that at
least five speakers make substantial contributions to each held-out partition and
that no individual voice can determine the evaluation result, while remaining
achievable given how unevenly speakers contribute to WAXAL.

\paragraph{Search procedure.}
The number of possible ways to divide 161 speakers into partitions of 113, 24, and
24 speakers is astronomically large, so evaluating every assignment was
infeasible. We therefore used a seeded randomized search. The speaker identifiers
were first sorted alphabetically, so that the search did not depend on the order in
which records happened to be read. In each of 50,000 trials, the sorted list was
randomly permuted using a pseudo-random number generator initialized with seed
42; the first 24 speakers in the permutation were assigned to validation, the next
24 to test, and the remaining 113 to training. The trial was discarded if it
violated the concentration constraint; otherwise $J(a)$ was computed and the
assignment was retained if it had the lowest value found so far. Using a fixed
seed made the search exactly reproducible: running the script again produces the
same partitions.

\paragraph{Version history.}
A first version of the procedure (version~1) evaluated 20,000 trials and assigned
16 speakers (10\% of 161) to each held-out partition, without a concentration
constraint. It achieved zero speaker overlap and well-balanced sizes, but one
speaker contributed 51.1\% of the validation utterances and another 35.1\% of the
test utterances. Version~1 was therefore retained only as a diagnostic pilot.
Version~2 added the 20\% concentration constraint, increased the number of
held-out speakers to 24 per partition, because with only 16 speakers the
constraint could rarely be satisfied, and increased the search to 50,000 trials,
because the constraint rejects many candidates. Version~2 was used for all
controlled experiments.

\paragraph{Result and verification.}
The selected version~2 assignment achieved $J(a) = 0.0029$. It contained 13,807
training utterances (79.76 hours, 113 speakers), 1,715 validation utterances
(9.59 hours, 24 speakers), and 1,671 test utterances (9.56 hours, 24 speakers).
Relative to the targets, the training partition was 0.4\% larger in utterances
and 0.8\% larger in duration, the validation partition was 0.3\% smaller in
utterances and 3.0\% smaller in duration, and the test partition was 2.8\% smaller
in utterances and 3.4\% smaller in duration. The largest single-speaker share was
19.5\% of utterances (19.8\% of duration) in validation and 17.2\% of utterances
(17.3\% of duration) in test. Automated assertions verified that every pooled
utterance was assigned exactly once, that utterance identifiers were unique, that
every audio file existed, that no transcription was empty, and that no speaker
appeared in more than one partition. The source and output manifests, the
speaker assignments, the constraints, the objective value, and all file hashes
were stored in a protocol summary file.

\subsubsection{Acoustic measurements for difficulty scoring}
\label{sec:snr}

Several curricula (Section~\ref{sec:curriculum}) rank utterances by their
signal-to-noise ratio (SNR), the ratio of the power of the speech to the power of
the background noise, expressed in decibels. A high SNR indicates clean speech
and a low SNR indicates noisy speech. WAXAL does not provide separate clean speech
and noise signals, so the true SNR cannot be computed; instead, we estimated an
SNR \emph{proxy} from the frame energies of each recording, decoded at its native
sampling rate of 48~kHz:

\begin{equation}
    \widehat{\mathrm{SNR}}_i = \min\!\left(60,\;
    20\log_{10}\frac{P_{75}\big(\{\operatorname{RMS}_f : f~\text{active}\}\big)}
    {P_{20}\big(\{\operatorname{RMS}_f : \text{all}~f\}\big)}\right)~\text{dB},
    \label{eq:snr-proxy}
\end{equation}

where $P_{75}$ denotes the 75th percentile and $P_{20}$ the 20th percentile, and
active frames are defined by Equation~\eqref{eq:active}. The 75th percentile of
the active frames represents a typical level of speech while being robust to a
few unusually loud frames, and the 20th percentile of all frames represents the
quiet background between words. The estimate was capped at 60~dB because, in
recordings with near-digital silence between words, the denominator approaches
zero and the ratio becomes arbitrarily large without reflecting any meaningful
difference in perceived quality. In the training partition the proxy ranged from
7.7 to 60.0~dB, with a median of 33.4~dB; 2,665 training utterances (19.3\%)
reached the 60~dB cap. The measurements were computed once, aligned to the
speaker-disjoint manifests by utterance identifier, and frozen with a SHA-256 hash
before any curriculum experiment was run.

\subsubsection{FLEURS out-of-domain evaluation data}
\label{sec:fleurs}

To evaluate out-of-domain generalization (RQ2), we used the Shona configuration
(\texttt{sn\_zw}) of FLEURS (Few-shot Learning Evaluation of Universal
Representations of Speech) \citep{Conneau2022}. FLEURS contains 2,463 training,
393 validation, and 925 test utterances for Shona. We chose FLEURS because it is,
to our knowledge, the only other publicly available, openly licensed corpus of
transcribed Shona speech with standard partitions, and because it differs from
WAXAL in every major respect that defines a domain. FLEURS consists of \emph{read}
speech: speakers read aloud sentences translated from English Wikipedia articles,
whereas WAXAL speakers spontaneously described images. FLEURS was therefore
recorded by different speakers, with different devices and recording conditions,
in a different speaking style, and on different topics and vocabulary. A model
that improves on FLEURS as well as on WAXAL is therefore likely to have learned
something general about Shona speech rather than the particularities of the WAXAL
collection. Evaluating ASR models on corpora other than the one used for training
is recommended practice for assessing robustness \citep{Likhomanenko2021}.

No FLEURS utterance was used for training, for selecting a curriculum, or for
tuning any hyperparameter. FLEURS results were computed only after the curriculum
choice had been fixed using WAXAL validation data, so that the out-of-domain
evaluation remained an independent check.

We applied a text contract with the same purpose and steps as for WAXAL
(Section~\ref{sec:normalization}): Unicode NFC normalization, case folding,
standardization of apostrophes and dashes, retention of letters, numbers,
combining marks, and apostrophes, replacement of other punctuation and symbols
with spaces, and collapsing of repeated spaces. The original transcription was
retained separately. The two implementations differ in minor details: the
FLEURS version selects characters by Unicode category (and therefore also keeps
separate combining marks and numeric forms such as fractions), whereas the
WAXAL version uses letter-and-digit tests and additionally maps backticks to
apostrophes. These differences affect very few characters. Every model was
scored against the same frozen FLEURS references, so comparisons between
models on FLEURS are unaffected, but absolute WAXAL and FLEURS WER values should
not be compared as if the scoring rules were identical.

An initial preparation of FLEURS removed digits as well as punctuation. Because
this deleted spoken numbers from 531 training transcriptions and created
mismatches between audio and text, runs that used that preparation were
excluded from all controlled comparisons. The initial manifests also reused the
same sentence identifier for several recordings of the same sentence, so the
corrected version assigned each recording a unique identifier derived from its
audio, while retaining the original sentence identifier for analysis.

After normalization, all 3,781 FLEURS utterances were retained in their original
partitions. An audit verified that every file was 16~kHz single-channel (mono)
audio as Whisper requires, that every manifest duration matched its audio, that
every utterance identifier was unique, and that no recording was duplicated
within FLEURS or shared with WAXAL (using the PCM hashes described in
Section~\ref{sec:leakage}). Because the hash includes the sampling rate, and
WAXAL is stored at 48~kHz while FLEURS is stored at 16~kHz, this cross-corpus
check can detect only exact copies, not the same recording stored at a different
rate. Nor can it establish whether FLEURS was part of Whisper's pretraining
data. FLEURS is therefore described as external to this study's fine-tuning
data, not as speech that Whisper has never encountered.

\subsubsection{Summary of data protocols}

Table~\ref{tab:protocols} summarizes the data versions used in the study. Only
the speaker-disjoint version~2 protocol was used for the controlled RQ1 and RQ2
experiments; the other versions are listed because they underlie exploratory
results referred to in this thesis.

\begin{table}[H]
    \centering
    \begin{tabular}{|p{4.2cm}|p{4.2cm}|r|r|r|}
        \hline
        Protocol & Purpose & Train & Validation & Test \\
        \hline
        Normalized original partitions & Exploratory baselines and model-size study & 13,799 & 1,683 & 1,711 \\
        \hline
        Speaker-disjoint version 1 & Diagnostic pilot only & 13,778 & 1,696 & 1,719 \\
        \hline
        Speaker-disjoint version 2 & All controlled experiments & 13,807 & 1,715 & 1,671 \\
        \hline
        FLEURS Shona (corrected) & Out-of-domain evaluation only & 2,463 & 393 & 925 \\
        \hline
    \end{tabular}
    \caption{Data protocols used in the study (number of utterances).}
    \label{tab:protocols}
\end{table}

\subsection{Model Input Preprocessing}
\label{sec:input-preprocessing}

Before training, the pipeline verified that every referenced audio file existed
and that every transcription was non-empty, so that a missing file could not
cause training to fail or silently skip examples part-way through.

\paragraph{Resampling.}
All audio was decoded as single-channel speech and resampled to 16~kHz, meaning
that the waveform was represented by 16,000 amplitude values per second. The
WAXAL recordings were stored at 48~kHz. Resampling to 16~kHz was required because
Whisper was pretrained exclusively on 16~kHz audio and its feature extractor
assumes that rate; supplying audio at any other rate would shift every frequency
the model sees. A 16~kHz rate represents frequencies up to 8~kHz (half the
sampling rate, by the sampling theorem \citep{Shannon1949}), which contains the
information most important for recognizing speech.

\paragraph{Log-Mel spectrogram.}
The Whisper feature extractor converted each waveform into an 80-channel
log-Mel spectrogram. A spectrogram is a time--frequency representation of sound:
the waveform is divided into short overlapping windows, and the energy at each
frequency is computed for each window using the short-time Fourier transform
(STFT). Whisper uses 25-millisecond windows (400 samples) advanced in steps of
10~milliseconds (160 samples), producing 100 frames per second
\citep{Radford2022}. The frequency axis is then mapped onto the Mel scale, a
perceptual scale on which equal distances correspond to equal perceived
differences in pitch \citep{Stevens1937}. It has finer resolution at low
frequencies, where most speech information lies, and coarser resolution at high
frequencies. The energies are summed into 80 Mel-spaced frequency bands and their
logarithm is taken, because human loudness perception is approximately
logarithmic and the logarithm compresses the very large range of speech energies
into a range suitable for a neural network \citep{Davis1980}. The result is a
matrix with 80 rows (frequency bands) and one column per 10-millisecond frame.

\paragraph{Fixed 30-second window.}
Whisper's encoder accepts exactly 3,000 frames, corresponding to 30 seconds of
audio, because it was pretrained on 30-second segments. Shorter utterances were
padded with silence-equivalent values to 3,000 frames. Utterances longer than 30
seconds were truncated to the first 30 seconds. In the speaker-disjoint protocol,
520 training utterances (3.8\%), 21 validation utterances, and 40 test utterances
exceeded 30 seconds; for these utterances the end of the audio was not seen by the
model while the reference transcription remained complete. This affects all
experimental conditions equally and therefore does not bias comparisons between
them, but it does introduce a small amount of label noise in training and
unavoidable deletion errors in evaluation, which are discussed as a limitation.

\paragraph{Attention mask.}
An attention mask, a binary vector marking which frames contain real audio and
which contain padding, was retained for each input. Whisper's encoder itself
always processes the complete 3,000-frame window, as during pretraining. The mask
was retained so that SpecAugment time masks (Section~\ref{sec:specaugment}) were
placed only within the real audio; a mask applied to padding would have no
effect and would silently weaken the augmentation for short utterances.

\paragraph{Tokenization.}
Target transcriptions were converted into sequences of integer tokens using the
Whisper multilingual tokenizer. A tokenizer splits text into sub-word units
drawn from a fixed vocabulary; Whisper uses a byte-level byte-pair encoding (BPE)
vocabulary \citep{Sennrich2016, Radford2022}, in which frequent character
sequences are represented by single tokens and rare words are split into several
tokens. The same vocabulary is shared across all of Whisper's 99 languages,
including Shona. The tokenizer was configured with the language set to Shona and
the task set to transcription. With this configuration, every target sequence
begins with a fixed prefix of special tokens,
\texttt{<|startoftranscript|><|sn|><|transcribe|><|notimestamps|>}, which tells
the decoder that the audio is Shona (\texttt{sn} is the language code for Shona),
that it should transcribe rather than translate into English, and that it should
not predict timestamps, and every sequence ends with an end-of-text token. Using
the same prefix in training and decoding ensured that the model was fine-tuned
for exactly the task on which it was evaluated.

\paragraph{Padding and the loss function.}
Training processes several utterances simultaneously in a \emph{batch}, which
must be stored as a rectangular array. Because transcriptions differ in length,
shorter token sequences were padded to the length of the longest sequence in the
batch. The padded positions do not correspond to real text, so the model must not
be trained to predict them. Each padded position was therefore replaced by the
value $-100$. The model was trained by minimizing the token-level cross-entropy
loss, which measures how much probability the model assigns to each correct
token:

\begin{equation}
    \mathcal{L} = -\frac{1}{|\mathcal{V}|}\sum_{(b,t)\in\mathcal{V}}
    \log p_\theta\!\left(y_{b,t} \mid y_{b,<t}, x_b\right),
    \qquad
    \mathcal{V} = \{(b,t) : y_{b,t} \neq -100\},
    \label{eq:ce-loss}
\end{equation}

where $b$ indexes the utterances in the batch, $t$ indexes token positions,
$y_{b,t}$ is the reference token, $y_{b,<t}$ are the preceding reference tokens,
$x_b$ is the input spectrogram, and $p_\theta$ is the probability predicted by the
model with parameters $\theta$. In PyTorch, the cross-entropy function excludes
every position whose target equals its \texttt{ignore\_index} parameter, whose
default value is $-100$. Positions labelled $-100$ are therefore excluded from both
the sum and the count $|\mathcal{V}|$: they contribute neither to the loss nor to
its gradient, and the loss is averaged over real tokens only. The leading
\texttt{<|startoftranscript|>} token was removed from the labels because the
Whisper decoder inserts it automatically as its first input.

\subsection{Baseline ASR Model}
\label{sec:baseline}

\subsubsection{Whisper}

Whisper \citep{Radford2022} is a family of ASR models released by OpenAI. Each
model is an encoder--decoder Transformer \citep{Vaswani2017}. The encoder converts
the log-Mel spectrogram into a sequence of internal representations of the audio,
and the decoder generates the transcription one token at a time, each time
conditioning on the encoder representations and on the tokens it has already
generated. The models were trained on 680,000 hours of audio paired with
transcripts collected from the internet. Of these, 438,000 hours were English
speech with English transcripts, 126,000 hours were non-English speech paired with
English translations, and 117,000 hours were non-English speech with transcripts
in the same language, covering 96 languages other than English. Because the
training data are dominated by English and a few high-resource languages, Whisper
performs considerably better on those languages than on low-resource languages
such as Shona, for which only a very small amount of pretraining data was
available. This gap is the motivation for fine-tuning Whisper on Shona.

Whisper was chosen as the foundation model for three reasons. First, Shona is one
of its supported languages, with its own language token, so the model could be
fine-tuned without modifying its vocabulary or architecture. Second, its
multilingual pretraining provides acoustic and linguistic knowledge that can be
transferred to a language with limited data. Third, its openly available
checkpoints in several sizes made it possible to study the effect of model
capacity under otherwise identical conditions.

\subsubsection{Model sizes}

During the exploratory phase we examined the Tiny, Base, Small, Medium, Large,
and Large-v3 checkpoints. The controlled experiments used Whisper Base
(74 million parameters; 6 encoder and 6 decoder layers with 512-dimensional
representations) as the screening model, and Whisper Medium (769 million
parameters; 24 encoder and 24 decoder layers with 1,024-dimensional
representations) for the model-capacity extension (Section~\ref{sec:medium}).

Whisper Base was chosen as the screening model because the RQ1 and RQ2 designs
required many conditions, each with three seeds; one Base training run took
approximately 0.6 GPU-hours, which made the complete design feasible. Whisper
Large was not carried into the controlled experiments for three reasons. First,
in the exploratory comparison on the original partitions, Large reduced test WER by
only a further 1.0 point relative to Medium (22.75\% versus 23.77\%), whereas
Medium reduced test WER by 4.7 points relative to Base. Second, one Large run took
15.35 hours, approximately six times as long as Medium (2.56 hours), so the full
design would have required several hundred GPU-hours. Third, Large did not fit in
24~GB of GPU memory with the training settings used for Base; it required a batch
size of one with gradient accumulation and an 8-bit optimizer, so a Large run
would have differed from the Base and Medium runs in its optimization procedure,
not only in its size, and any difference could not have been attributed to model
capacity alone. Tiny and Small were not carried forward because Base and Medium
already bracketed the range of capacities that was computationally feasible.

\subsubsection{Zero-shot evaluation}

The pinned Whisper Base checkpoint was also evaluated \emph{zero-shot}, that is,
without any training on Shona, using the same speaker-disjoint partitions,
FLEURS references, language and task prefix, decoding settings, and text
contract as its fine-tuned counterparts. This established how much of the final
performance was due to fine-tuning rather than to the pretrained model. The
zero-shot evaluation was added after the fine-tuning experiments rather than
before them, which does not affect its validity because it involves no training
or selection. A zero-shot evaluation of Whisper Medium had not been completed at
the time of writing. \todo{Update if the Medium zero-shot evaluation is run.}
``Zero-shot'' here means that no Shona fine-tuning was performed in this study;
it does not imply that Whisper's pretraining data contained no Shona. Because WER counts insertions as errors, it can
exceed 100\% when a model produces many words that are not in the reference, as
an untrained model typically does.

\subsubsection{Fine-tuning configuration}
\label{sec:finetuning}

Every controlled run started from the pinned pretrained checkpoint (Whisper Base
revision \texttt{e37978b9} and Whisper Medium revision \texttt{abdf7c39}); no run
was initialized from the weights of a previous experiment, so that each condition
was independent of every other. All parameters of the model were updated
(full fine-tuning). The following settings were identical for every controlled
run, for both model sizes.

\begin{itemize}
    \item \textbf{Optimizer:} AdamW \citep{Loshchilov2019}, the default optimizer
    of the Transformers training loop. AdamW adapts the step size of each
    parameter according to running averages of its past gradients and squared
    gradients, which makes training robust to differences in scale between
    parameters. The pinned library defaults were used: moment coefficients
    $\beta_1 = 0.9$ and $\beta_2 = 0.999$, numerical-stability constant
    $\epsilon = 10^{-8}$, and gradients clipped to a maximum norm of 1.0. Weight
    decay was zero, so although AdamW supports decoupled weight decay, no weight
    penalty was applied in these runs. Label smoothing was also zero. These
    inherited defaults were not tuned and were identical across conditions.
    \item \textbf{Learning rate:} $1 \times 10^{-5}$. The learning rate controls
    the size of each parameter update. A small learning rate was chosen because
    large updates can overwrite the knowledge acquired during pretraining, a
    problem known as catastrophic forgetting; small learning rates are standard
    for fine-tuning large pretrained Transformers \citep{Mosbach2021}, and
    $1 \times 10^{-5}$ is the value used in the widely adopted Whisper fine-tuning
    recipe \citep{Gandhi2022}.
    \item \textbf{Warm-up:} 500 steps, during which the learning rate increased
    linearly from zero to $1 \times 10^{-5}$, after which it decreased linearly
    towards zero by the end of training. At the start of fine-tuning, the
    optimizer's running estimates of gradient statistics are unreliable and the
    model has not yet adapted to the new data, so full-size updates can be
    unstable; a gradual warm-up avoids large, damaging early updates
    \citep{Goyal2017, Mosbach2021}. The value of 500 steps, approximately 7\% of
    the 6,906 total updates, follows the same recipe \citep{Gandhi2022}.
    \item \textbf{Batch size:} six utterances per update, with no gradient
    accumulation. This was the largest batch that fitted comfortably in GPU memory
    for Base. The same batch size was retained for Medium, which fitted with a
    peak memory reservation of approximately 15~GB, so that both model sizes
    received exactly the same sequence of updates.
    \item \textbf{Epochs:} three complete passes over the training data. With
    13,807 utterances and a batch size of six, one epoch comprises 2,302 updates
    (the final batch of each epoch contains one utterance), and three epochs
    comprise 6,906 updates. Three epochs was the largest number that allowed every
    condition to be trained with three seeds within the available computing budget.
    The training budget was fixed in advance and identical for all conditions,
    because curricula are known to change how quickly a model learns; comparing
    conditions at the same number of updates is therefore essential for a fair
    comparison. Conclusions are consequently conditional on this budget.
    \item \textbf{Mixed precision (FP16):} most calculations were performed in
    16-bit floating-point format rather than 32-bit. This approximately halves the
    memory needed for activations and substantially increases speed on modern
    GPUs, while a master copy of the weights in 32-bit precision and dynamic loss
    scaling preserve numerical accuracy \citep{Micikevicius2018}. Dynamic loss
    scaling multiplies the loss by a large factor so that small gradients do not
    round to zero in 16-bit format; if the scaled gradients overflow, the update is
    skipped and the factor is reduced. A small number of such skipped updates
    occurred at the very start of each run.
    \item \textbf{Gradient checkpointing:} instead of storing every intermediate
    result during the forward pass for use in the backward pass, only selected
    results were stored and the rest were recomputed when needed. This reduces
    memory use substantially at the cost of approximately one additional forward
    computation \citep{Chen2016}, and was what allowed Medium to be trained with
    the same batch size as Base on a 24~GB GPU.
    \item \textbf{Deterministic execution:} deterministic GPU algorithms were
    enabled so that re-running a configuration on the same hardware and software
    reproduced the same model; this was verified by an exact re-run of one
    dynamic-curriculum condition (C5, seed 42), which reproduced the checkpoint
    hashes and metrics. Determinism is not guaranteed across different GPU types
    or library versions.
\end{itemize}

\subsubsection{Decoding}
\label{sec:decoding}

Transcriptions were generated with greedy decoding, in which the decoder selects
the single most probable token at each step. Greedy decoding was used for all
controlled experiments because it is deterministic, so that differences between
conditions cannot arise from randomness in decoding, and because it is
considerably cheaper than beam search, which tracks several candidate
transcriptions in parallel. Using one fixed decoding method also ensured that
decoding did not become an additional factor in the comparison of training
strategies; we do not claim that greedy decoding gives the best achievable WER.
The maximum output length was 225 tokens, which is sufficient for a 30-second
Shona utterance and prevents a model that has begun repeating itself from
generating indefinitely. The language and task were set explicitly to Shona
transcription. The legacy forced-decoder-token setting and the default list of
suppressed tokens were cleared in every run, following the standard Whisper
fine-tuning recipe \citep{Gandhi2022}, so that the language and task were
specified in only one place and every condition decoded with exactly the same
vocabulary.

\subsubsection{Checkpointing and model selection}
\label{sec:checkpointing}

The model was evaluated on the validation partition, and a checkpoint (a saved
copy of the model) was written, at the end of each epoch, that is, after updates
2,302, 4,604, and 6,906. Evaluating at epoch boundaries, rather than every fixed
number of updates, ensured that every saved model had seen each training
utterance the same number of times. This matters because, part-way through an
epoch, a curriculum model has seen only the easiest utterances whereas a randomly
ordered model has seen a representative sample, so mid-epoch checkpoints would
compare models that had been exposed to systematically different data. The
checkpoint with the lowest validation WER was selected as the final model and
evaluated once on the test partition. The test partition was never used to select
checkpoints or hyperparameters.

% Input Part 2: Curriculum, Augmentation, Medium Extension, and Evaluation Metrics
\input{chapter3_methodology_part2.tex}

\begin{thebibliography}{99}

\bibitem[Bakhturina et~al.(2021)]{Bakhturina2021}
Bakhturina, E., Lavrukhin, V., and Ginsburg, B. (2021).
A toolbox for construction and analysis of speech datasets.
\emph{NeurIPS Datasets and Benchmarks}.
\url{https://arxiv.org/abs/2104.04896}.

\bibitem[Bengio et~al.(2009)]{Bengio2009}
Bengio, Y., Louradour, J., Collobert, R., and Weston, J. (2009).
Curriculum learning. \emph{Proceedings of ICML}.
\url{https://doi.org/10.1145/1553374.1553380}.

\bibitem[Biewald(2020)]{Biewald2020}
Biewald, L. (2020). Experiment tracking with Weights and Biases.
Software. \url{https://wandb.ai/site/research/}.

\bibitem[Chen et~al.(2016)]{Chen2016}
Chen, T., Xu, B., Zhang, C., and Guestrin, C. (2016).
Training deep nets with sublinear memory cost.
\url{https://arxiv.org/abs/1604.06174}.

\bibitem[Conneau et~al.(2022)]{Conneau2022}
Conneau, A., et~al. (2022). FLEURS: Few-shot learning evaluation of universal
representations of speech. \url{https://arxiv.org/abs/2205.12446}.

\bibitem[Davis and Mermelstein(1980)]{Davis1980}
Davis, S.~B., and Mermelstein, P. (1980).
Comparison of parametric representations for monosyllabic word recognition
in continuously spoken sentences. \emph{IEEE Transactions on Acoustics,
Speech, and Signal Processing}, 28(4), 357--366.
\url{https://doi.org/10.1109/TASSP.1980.1163420}.

\bibitem[Dettmers et~al.(2022)]{Dettmers2022}
Dettmers, T., Lewis, M., Shleifer, S., and Zettlemoyer, L. (2022).
8-bit optimizers via block-wise quantization. \emph{Proceedings of ICLR}.
\url{https://arxiv.org/abs/2110.02861}.

\bibitem[Diack et~al.(2026)]{Diack2026}
Diack, A., et~al. (2026). WAXAL: A large-scale multilingual African language
speech corpus. arXiv:2602.02734, version 3.
\url{https://arxiv.org/abs/2602.02734}.

\bibitem[Dodge et~al.(2020)]{Dodge2020}
Dodge, J., Ilharco, G., Schwartz, R., Farhadi, A., Hajishirzi, H., and
Smith, N.~A. (2020). Fine-tuning pretrained language models: Weight
initializations, data orders, and early stopping.
\url{https://arxiv.org/abs/2002.06305}.

\bibitem[Gandhi(2022)]{Gandhi2022}
Gandhi, S. (2022). Fine-tune Whisper for multilingual ASR with Transformers.
\emph{Hugging Face Blog}.
\url{https://huggingface.co/blog/fine-tune-whisper}.

\bibitem[Garofolo et~al.(1993)]{Garofolo1993}
Garofolo, J.~S., Lamel, L.~F., Fisher, W.~M., Fiscus, J.~G., Pallett, D.~S.,
Dahlgren, N.~L., and Zue, V. (1993). TIMIT acoustic-phonetic continuous speech
corpus. \emph{Linguistic Data Consortium}, LDC93S1.
\url{https://catalog.ldc.upenn.edu/LDC93S1}.

\bibitem[Goyal et~al.(2017)]{Goyal2017}
Goyal, P., et~al. (2017). Accurate, large minibatch SGD: Training ImageNet
in 1 hour. \url{https://arxiv.org/abs/1706.02677}.

\bibitem[Kaufman et~al.(2012)]{Kaufman2012}
Kaufman, S., Rosset, S., Perlich, C., and Stitelman, O. (2012).
Leakage in data mining: Formulation, detection, and avoidance.
\emph{ACM Transactions on Knowledge Discovery from Data}, 6(4).
\url{https://doi.org/10.1145/2382577.2382579}.

\bibitem[Koenecke et~al.(2024)]{Koenecke2024}
Koenecke, A., Choi, A.~S.~G., Mei, K.~X., Schellmann, H., and Sloane, M.
(2024). Careless Whisper: Speech-to-text hallucination harms.
\url{https://arxiv.org/abs/2402.08021}.

\bibitem[Lhoest et~al.(2021)]{Lhoest2021}
Lhoest, Q., et~al. (2021). Datasets: A community library for natural language
processing. \emph{Proceedings of EMNLP: System Demonstrations}, 175--184.
\url{https://aclanthology.org/2021.emnlp-demo.21/}.

\bibitem[Likhomanenko et~al.(2021)]{Likhomanenko2021}
Likhomanenko, T., et~al. (2021).
Rethinking evaluation in ASR: Are our models robust enough?
\emph{Proceedings of Interspeech}, 311--315.
\url{https://doi.org/10.21437/Interspeech.2021-1758}.

\bibitem[Loshchilov and Hutter(2019)]{Loshchilov2019}
Loshchilov, I., and Hutter, F. (2019). Decoupled weight decay regularization.
\emph{Proceedings of ICLR}. \url{https://arxiv.org/abs/1711.05101}.

\bibitem[McFee et~al.(2015)]{McFee2015}
McFee, B., et~al. (2015). librosa: Audio and music signal analysis in Python.
\emph{Proceedings of the Python in Science Conference}.
\url{https://doi.org/10.25080/Majora-7b98e3ed-003}.

\bibitem[Micikevicius et~al.(2018)]{Micikevicius2018}
Micikevicius, P., et~al. (2018). Mixed precision training.
\emph{Proceedings of ICLR}. \url{https://arxiv.org/abs/1710.03740}.

\bibitem[Mosbach et~al.(2021)]{Mosbach2021}
Mosbach, M., Andriushchenko, M., and Klakow, D. (2021).
On the stability of fine-tuning BERT: Misconceptions, explanations, and strong
baselines. \emph{Proceedings of ICLR}.
\url{https://arxiv.org/abs/2006.04884}.

\bibitem[Panayotov et~al.(2015)]{Panayotov2015}
Panayotov, V., Chen, G., Povey, D., and Khudanpur, S. (2015).
LibriSpeech: An ASR corpus based on public domain audio books.
\emph{Proceedings of ICASSP}.
\url{https://doi.org/10.1109/ICASSP.2015.7178964}.

\bibitem[Paszke et~al.(2019)]{Paszke2019}
Paszke, A., et~al. (2019). PyTorch: An imperative style, high-performance deep
learning library. \emph{Advances in Neural Information Processing Systems}, 32.
\url{https://arxiv.org/abs/1912.01703}.

\bibitem[Radford et~al.(2022)]{Radford2022}
Radford, A., Kim, J.~W., Xu, T., Brockman, G., McLeavey, C., and Sutskever, I.
(2022). Robust speech recognition via large-scale weak supervision.
arXiv:2212.04356. \url{https://arxiv.org/abs/2212.04356}.

\bibitem[Reimers and Gurevych(2017)]{Reimers2017}
Reimers, N., and Gurevych, I. (2017). Reporting score distributions makes a
difference: Performance study of LSTM-networks for sequence tagging.
\emph{Proceedings of EMNLP}.
\url{https://doi.org/10.18653/v1/D17-1035}.

\bibitem[Sennrich et~al.(2016)]{Sennrich2016}
Sennrich, R., Haddow, B., and Birch, A. (2016).
Neural machine translation of rare words with subword units.
\emph{Proceedings of ACL}, 1715--1725.
\url{https://aclanthology.org/P16-1162/}.

\bibitem[Shannon(1949)]{Shannon1949}
Shannon, C.~E. (1949). Communication in the presence of noise.
\emph{Proceedings of the IRE}, 37(1), 10--21.
\url{https://doi.org/10.1109/JRPROC.1949.232969}.

\bibitem[Stevens et~al.(1937)]{Stevens1937}
Stevens, S.~S., Volkmann, J., and Newman, E.~B. (1937).
A scale for the measurement of the psychological magnitude pitch.
\emph{The Journal of the Acoustical Society of America}, 8(3), 185--190.
\url{https://doi.org/10.1121/1.1915893}.

\bibitem[Unicode Consortium(n.d.)]{UAX15}
Unicode Consortium (n.d.). Unicode Standard Annex \#15: Unicode normalization
forms. \url{https://www.unicode.org/reports/tr15/}.

\bibitem[Vaswani et~al.(2017)]{Vaswani2017}
Vaswani, A., et~al. (2017). Attention is all you need.
\emph{Advances in Neural Information Processing Systems}, 30.
\url{https://arxiv.org/abs/1706.03762}.

\bibitem[Wolf et~al.(2020)]{Wolf2020}
Wolf, T., et~al. (2020). Transformers: State-of-the-art natural language
processing. \emph{Proceedings of EMNLP: System Demonstrations}, 38--45.
\url{https://aclanthology.org/2020.emnlp-demos.6/}.

\bibitem[Yuan et~al.(2006)]{Yuan2006}
Yuan, J., Liberman, M., and Cieri, C. (2006).
Towards an integrated understanding of speaking rate in conversation.
\emph{Proceedings of Interspeech}.
\url{https://www.isca-archive.org/interspeech_2006/yuan06_interspeech.html}.

\end{thebibliography}

\end{document}
